
\documentclass[10pt]{article}
\usepackage[utf8]{inputenc}
\usepackage[T1]{fontenc}
\usepackage{amsmath, amssymb, xcolor, geometry, enumitem, multicol, tcolorbox}

\definecolor{mainblue}{RGB}{0, 102, 204}
\geometry{a4paper, margin=1.2cm, top=1cm, bottom=3cm, footskip=1.2cm}

\begin{document}
\enlargethispage{2.5cm}

% Modern Header
\begin{tcolorbox}[colback=mainblue!5, colframe=mainblue, sharp corners, boxrule=0.5pt]
    \small \textbf{Unit:} Ὅ6 Unit 7: Factoring \hfill \textbf{Name:} \rule{3cm}{0.4pt} \\
    \small \textbf{Topic:} Expanding Special Patterns & Perfect Square Trinomial \hfill \textbf{Date:} \rule{2cm}{0.4pt} \\
    \centering \textbf{\large Challenge (Habanero)}
\end{tcolorbox}

\vspace{0.2cm}

% MCQ Section
\begin{multicols}{2}
\begin{enumerate}[label=\textbf{Q\arabic*.}, leftmargin=*, itemsep=6pt, parsep=0pt]

  \item \small What is the expansion of $(x + 3)^2$?
  \begin{enumerate}[label=(\Alph*), leftmargin=0.6cm, nosep, itemsep=2pt]
    \item $x^2 + 6x + 9$
    \item $x^2 - 6x + 9$
    \item $x^2 + 3x + 9$
    \item $x^2 - 3x + 9$
    \item $x^2 + 9$
  \end{enumerate}

  \item \small A gaming company uses the expression $x^2 - 10x + 25$ to model the score of a new game. What is the factored form of this expression?
  \begin{enumerate}[label=(\Alph*), leftmargin=0.6cm, nosep, itemsep=2pt]
    \item $(x - 5)^2$
    \item $(x + 5)^2$
    \item $(x - 2)(x - 5)$
    \item $(x + 2)(x + 5)$
    \item $(x - 1)(x - 10)$
  \end{enumerate}

  \item \small A coder is debugging a line of code that contains the expression $(2x + 5)^2$. What is the expanded form of this expression?
  \begin{enumerate}[label=(\Alph*), leftmargin=0.6cm, nosep, itemsep=2pt]
    \item $4x^2 + 20x + 25$
    \item $4x^2 - 20x + 25$
    \item $2x^2 + 10x + 25$
    \item $2x^2 - 10x + 25$
    \item $4x^2 + 10x + 25$
  \end{enumerate}

  \item \small A data analyst is working with the expression $9x^2 - 12x + 4$. What type of trinomial is this?
  \begin{enumerate}[label=(\Alph*), leftmargin=0.6cm, nosep, itemsep=2pt]
    \item Perfect square trinomial
    \item Difference of squares
    \item Sum of cubes
    \item Difference of cubes
    \item None of the above
  \end{enumerate}

  \item \small What is the factored form of the expression $16x^2 - 24x + 9$?
  \begin{enumerate}[label=(\Alph*), leftmargin=0.6cm, nosep, itemsep=2pt]
    \item $(4x - 3)^2$
    \item $(4x + 3)^2$
    \item $(2x - 3)(2x - 3)$
    \item $(2x + 3)(2x + 3)$
    \item $(4x - 1)(4x - 9)$
  \end{enumerate}

  \item \small A gaming platform uses the expression $(3x - 2)^2$ to model the number of players online. What is the expanded form of this expression?
  \begin{enumerate}[label=(\Alph*), leftmargin=0.6cm, nosep, itemsep=2pt]
    \item $9x^2 - 12x + 4$
    \item $9x^2 + 12x + 4$
    \item $9x^2 - 6x + 4$
    \item $9x^2 + 6x + 4$
    \item $9x^2 - 4x + 4$
  \end{enumerate}

  \item \small What is the expansion of $(2x + 1)^2$?
  \begin{enumerate}[label=(\Alph*), leftmargin=0.6cm, nosep, itemsep=2pt]
    \item $4x^2 + 4x + 1$
    \item $4x^2 - 4x + 1$
    \item $2x^2 + 4x + 1$
    \item $2x^2 - 4x + 1$
    \item $4x^2 + 1$
  \end{enumerate}

  \item \small A software company uses the expression $x^2 + 14x + 49$ to model the number of users. What is the factored form of this expression?
  \begin{enumerate}[label=(\Alph*), leftmargin=0.6cm, nosep, itemsep=2pt]
    \item $(x + 7)^2$
    \item $(x - 7)^2$
    \item $(x + 1)(x + 49)$
    \item $(x - 1)(x - 49)$
    \item $(x + 2)(x + 24)$
  \end{enumerate}

\end{enumerate}
\end{multicols}

\vspace{-0.2cm}
\hrule
\vspace{0.2cm}
\textbf{\small Open Ended Questions (FRQ)}
\begin{enumerate}[label=\textbf{Q\arabic*.}, start=9, itemsep=5pt, topsep=0pt]

  \item \small Factor the expression $25x^2 - 30x + 9$. Show all work and explain your reasoning. \vspace{2cm}

  \item \small Expand the expression $(x + 2)(x - 5)$ and then factor the result. Explain the relationship between the original expression and the factored form. \vspace{2cm}

\end{enumerate}

\vfill

% Chef's Tips Footer
\begin{tcolorbox}[colback=blue!5, colframe=mainblue!50, title=\small \textbf{Chef's Tips}, fonttitle=\bfseries, bottom=1mm]
\begin{itemize}[noitemsep, topsep=2pt, leftmargin=0.5cm]
  \item \scriptsize Remind students to simplify their answers\item \scriptsize Encourage students to check their work\item \scriptsize Emphasize the importance of showing all work and explaining reasoning
\end{itemize}
\end{tcolorbox}

\end{document}
  